%!TeX root=../gardentop.tex
\chapter{»It Has Come!«} 
	
\begin{figure}[t!]
\centering
\includegraphics[width=\linewidth]{september}
\end{figure}

 \lettrine[lines=6]{O}{f} course Dr~Craven had been sent for the morning after Colin had had his tantrum. He was always sent for at once when such a thing occurred and he always found, when he arrived, a white shaken boy lying on his bed, sulky and still so hysterical that he was ready to break into fresh sobbing at the least word. In fact, Dr~Craven dreaded and detested the difficulties of these visits. On this occasion he was away from Misselthwaite Manor until afternoon.

»How is he?« he asked Mrs~Medlock rather irritably when he arrived. »He will break a blood-vessel in one of those fits some day. The boy is half insane with hysteria and self-indulgence.«

»Well, sir,« answered Mrs~Medlock, »you'll scarcely believe your eyes when you see him. That plain sour-faced child that's almost as bad as himself has just bewitched him. How she's done it there's no telling. The Lord knows she's nothing to look at and you scarcely ever hear her speak, but she did what none of us dare do. She just flew at him like a little cat last night, and stamped her feet and ordered him to stop screaming, and somehow she startled him so that he actually did stop, and this afternoon—well just come up and see, sir. It's past crediting.«

The scene which Dr~Craven beheld when he entered his patient's room was indeed rather astonishing to him. As Mrs~Medlock opened the door he heard laughing and chattering. Colin was on his sofa in his dressing-gown and he was sitting up quite straight looking at a picture in one of the garden books and talking to the plain child who at that moment could scarcely be called plain at all because her face was so glowing with enjoyment.

»Those long spires of blue ones—we'll have a lot of those,« Colin was announcing. »They're called Del-phin-iums.«

»Dickon says they're larkspurs made big and grand,« cried Mistress Mary. »There are clumps there already.«

Then they saw Dr~Craven and stopped. Mary became quite still and Colin looked fretful.

»I am sorry to hear you were ill last night, my boy,« Dr~Craven said a trifle nervously. He was rather a nervous man.

»I'm better now—much better,« Colin answered, rather like a Rajah. »I'm going out in my chair in a day or two if it is fine. I want some fresh air.«

Dr~Craven sat down by him and felt his pulse and looked at him curiously.

»It must be a very fine day,« he said, »and you must be very careful not to tire yourself.«

»Fresh air won't tire me,« said the young Rajah.

As there had been occasions when this same young gentleman had shrieked aloud with rage and had insisted that fresh air would give him cold and kill him, it is not to be wondered at that his doctor felt somewhat startled.

»I thought you did not like fresh air,« he said.

»I don't when I am by myself,« replied the Rajah; »but my cousin is going out with me.«

»And the nurse, of course?« suggested Dr~Craven.

»No, I will not have the nurse,« so magnificently that Mary could not help remembering how the young native Prince had looked with his diamonds and emeralds and pearls stuck all over him and the great rubies on the small dark hand he had waved to command his servants to approach with salaams and receive his orders.

»My cousin knows how to take care of me. I am always better when she is with me. She made me better last night. A very strong boy I know will push my carriage.«

Dr~Craven felt rather alarmed. If this tiresome hysterical boy should chance to get well he himself would lose all chance of inheriting Misselthwaite; but he was not an unscrupulous man, though he was a weak one, and he did not intend to let him run into actual danger.

»He must be a strong boy and a steady boy,« he said. »And I must know something about him. Who is he? What is his name?«

»It's Dickon,« Mary spoke up suddenly. She felt somehow that everybody who knew the moor must know Dickon. And she was right, too. She saw that in a moment Dr~Craven's serious face relaxed into a relieved smile.

»Oh, Dickon,« he said. »If it is Dickon you will be safe enough. He's as strong as a moor pony, is Dickon.«

»And he's trusty,« said Mary. »He's th' trustiest lad i' Yorkshire.« She had been talking Yorkshire to Colin and she forgot herself.

»Did Dickon teach you that?« asked Dr~Craven, laughing outright.

»I'm learning it as if it was French,« said Mary rather coldly. »It's like a native dialect in India. Very clever people try to learn them. I like it and so does Colin.«

»Well, well,« he said. »If it amuses you perhaps it won't do you any harm. Did you take your bromide last night, Colin?«

»No,« Colin answered. »I wouldn't take it at first and after Mary made me quiet she talked me to sleep—in a low voice—about the spring creeping into a garden.«

»That sounds soothing,« said Dr~Craven, more perplexed than ever and glancing sideways at Mistress Mary sitting on her stool and looking down silently at the carpet. »You are evidently better, but you must remember—«

»I don't want to remember,« interrupted the Rajah, appearing again. »When I lie by myself and remember I begin to have pains everywhere and I think of things that make me begin to scream because I hate them so. If there was a doctor anywhere who could make you forget you were ill instead of remembering it I would have him brought here.« And he waved a thin hand which ought really to have been covered with royal signet rings made of rubies. »It is because my cousin makes me forget that she makes me better.«

Dr~Craven had never made such a short stay after a »tantrum«; usually he was obliged to remain a very long time and do a great many things. This afternoon he did not give any medicine or leave any new orders and he was spared any disagreeable scenes. When he went down-stairs he looked very thoughtful and when he talked to Mrs~Medlock in the library she felt that he was a much puzzled man.

»Well, sir,« she ventured, »could you have believed it?«

»It is certainly a new state of affairs,« said the doctor. »And there's no denying it is better than the old one.«

»I believe Susan Sowerby's right—I do that,« said Mrs~Medlock. »I stopped in her cottage on my way to Thwaite yesterday and had a bit of talk with her. And she says to me, »Well, Sarah Ann, she mayn't be a good child, an' she mayn't be a pretty one, but she's a child, an' children needs children.« We went to school together, Susan Sowerby and me.«

»She's the best sick nurse I know,« said Dr~Craven. »When I find her in a cottage I know the chances are that I shall save my patient.«

Mrs~Medlock smiled. She was fond of Susan Sowerby.

»She's got a way with her, has Susan,« she went on quite volubly. »I've been thinking all morning of one thing she said yesterday. She says, »Once when I was givin' th' children a bit of a preach after they'd been fightin' I ses to 'em all, »When I was at school my jography told as th' world was shaped like a orange an' I found out before I was ten that th' whole orange doesn't belong to nobody. No one owns more than his bit of a quarter an' there's times it seems like there's not enow quarters to go round. But don't you—none o' you—think as you own th' whole orange or you'll find out you're mistaken, an' you won't find it out without hard knocks.« What children learns from children,« she says, »is that there's no sense in grabbin' at th' whole orange—peel an' all. If you do you'll likely not get even th' pips, an' them's too bitter to eat.««

»She's a shrewd woman,« said Dr~Craven, putting on his coat.

»Well, she's got a way of saying things,« ended Mrs~Medlock, much pleased. »Sometimes I've said to her, 'Eh! Susan, if you was a different woman an' didn't talk such broad Yorkshire I've seen the times when I should have said you was clever.'«

That night Colin slept without once awakening and when he opened his eyes in the morning he lay still and smiled without knowing it—smiled because he felt so curiously comfortable. It was actually nice to be awake, and he turned over and stretched his limbs luxuriously. He felt as if tight strings which had held him had loosened themselves and let him go. He did not know that Dr~Craven would have said that his nerves had relaxed and rested themselves. Instead of lying and staring at the wall and wishing he had not awakened, his mind was full of the plans he and Mary had made yesterday, of pictures of the garden and of Dickon and his wild creatures. It was so nice to have things to think about. And he had not been awake more than ten minutes when he heard feet running along the corridor and Mary was at the door. The next minute she was in the room and had run across to his bed, bringing with her a waft of fresh air full of the scent of the morning.

»You've been out! You've been out! There's that nice smell of leaves!« he cried.

She had been running and her hair was loose and blown and she was bright with the air and pink-cheeked, though he could not see it.

»It's so beautiful!« she said, a little breathless with her speed. »You never saw anything so beautiful! It has come! I thought it had come that other morning, but it was only coming. It is here now! It has come, the Spring! Dickon says so!«

»Has it?« cried Colin, and though he really knew nothing about it he felt his heart beat. He actually sat up in bed.

»Open the window!« he added, laughing half with joyful excitement and half at his own fancy. »Perhaps we may hear golden trumpets!«

And though he laughed, Mary was at the window in a moment and in a moment more it was opened wide and freshness and softness and scents and birds' songs were pouring through.

»That's fresh air,« she said. »Lie on your back and draw in long breaths of it. That's what Dickon does when he's lying on the moor. He says he feels it in his veins and it makes him strong and he feels as if he could live forever and ever. Breathe it and breathe it.«

She was only repeating what Dickon had told her, but she caught Colin's fancy.

»»Forever and ever«! Does it make him feel like that?« he said, and he did as she told him, drawing in long deep breaths over and over again until he felt that something quite new and delightful was happening to him.

Mary was at his bedside again.

»Things are crowding up out of the earth,« she ran on in a hurry. »And there are flowers uncurling and buds on everything and the green veil has covered nearly all the gray and the birds are in such a hurry about their nests for fear they may be too late that some of them are even fighting for places in the secret garden. And the rose-bushes look as wick as wick can be, and there are primroses in the lanes and woods, and the seeds we planted are up, and Dickon has brought the fox and the crow and the squirrels and a new-born lamb.«

And then she paused for breath. The new-born lamb Dickon had found three days before lying by its dead mother among the gorse bushes on the moor. It was not the first motherless lamb he had found and he knew what to do with it. He had taken it to the cottage wrapped in his jacket and he had let it lie near the fire and had fed it with warm milk. It was a soft thing with a darling silly baby face and legs rather long for its body. Dickon had carried it over the moor in his arms and its feeding bottle was in his pocket with a squirrel, and when Mary had sat under a tree with its limp warmness huddled on her lap she had felt as if she were too full of strange joy to speak. A lamb—a lamb! A living lamb who lay on your lap like a baby!

She was describing it with great joy and Colin was listening and drawing in long breaths of air when the nurse entered. She started a little at the sight of the open window. She had sat stifling in the room many a warm day because her patient was sure that open windows gave people cold.

»Are you sure you are not chilly, Master Colin?« she inquired.

»No,« was the answer. »I am breathing long breaths of fresh air. It makes you strong. I am going to get up to the sofa for breakfast and my cousin will have breakfast with me.«

The nurse went away, concealing a smile, to give the order for two breakfasts. She found the servants' hall a more amusing place than the invalid's chamber and just now everybody wanted to hear the news from up-stairs. There was a great deal of joking about the unpopular young recluse who, as the cook said, »had found his master, and good for him.« The servants' hall had been very tired of the tantrums, and the butler, who was a man with a family, had more than once expressed his opinion that the invalid would be all the better »for a good hiding.«

When Colin was on his sofa and the breakfast for two was put upon the table he made an announcement to the nurse in his most Rajah-like manner.

»A boy, and a fox, and a crow, and two squirrels, and a new-born lamb, are coming to see me this morning. I want them brought up-stairs as soon as they come,« he said. »You are not to begin playing with the animals in the servants' hall and keep them there. I want them here.«

The nurse gave a slight gasp and tried to conceal it with a cough.

»Yes, sir,« she answered.

»I'll tell you what you can do,« added Colin, waving his hand. »You can tell Martha to bring them here. The boy is Martha's brother. His name is Dickon and he is an animal charmer.«

»I hope the animals won't bite, Master Colin,« said the nurse.

»I told you he was a charmer,« said Colin austerely. »Charmers' animals never bite.«

»There are snake-charmers in India,« said Mary; »and they can put their snakes' heads in their mouths.«

»Goodness!« shuddered the nurse.

They ate their breakfast with the morning air pouring in upon them. Colin's breakfast was a very good one and Mary watched him with serious interest.

»You will begin to get fatter just as I did,« she said. »I never wanted my breakfast when I was in India and now I always want it.«

»I wanted mine this morning,« said Colin. »Perhaps it was the fresh air. When do you think Dickon will come?«

He was not long in coming. In about ten minutes Mary held up her hand.

»Listen!« she said. »Did you hear a caw?«

Colin listened and heard it, the oddest sound in the world to hear inside a house, a hoarse »caw-caw.«

»Yes,« he answered.

»That's Soot,« said Mary. »Listen again! Do you hear a bleat—a tiny one?«

»Oh, yes!« cried Colin, quite flushing.

»That's the new-born lamb,« said Mary. »He's coming.«

Dickon's moorland boots were thick and clumsy and though he tried to walk quietly they made a clumping sound as he walked through the long corridors. Mary and Colin heard him marching—marching, until he passed through the tapestry door on to the soft carpet of Colin's own passage.

»If you please, sir,« announced Martha, opening the door, »if you please, sir, here's Dickon an' his creatures.«

Dickon came in smiling his nicest wide smile. The new-born lamb was in his arms and the little red fox trotted by his side. Nut sat on his left shoulder and Soot on his right and Shell's head and paws peeped out of his coat pocket.

Colin slowly sat up and stared and stared—as he had stared when he first saw Mary; but this was a stare of wonder and delight. The truth was that in spite of all he had heard he had not in the least understood what this boy would be like and that his fox and his crow and his squirrels and his lamb were so near to him and his friendliness that they seemed almost to be part of himself. Colin had never talked to a boy in his life and he was so overwhelmed by his own pleasure and curiosity that he did not even think of speaking.

But Dickon did not feel the least shy or awkward. He had not felt embarrassed because the crow had not known his language and had only stared and had not spoken to him the first time they met. Creatures were always like that until they found out about you. He walked over to Colin's sofa and put the new-born lamb quietly on his lap, and immediately the little creature turned to the warm velvet dressing-gown and began to nuzzle and nuzzle into its folds and butt its tight-curled head with soft impatience against his side. Of course no boy could have helped speaking then.

»What is it doing?« cried Colin. »What does it want?«

»It wants its mother,« said Dickon, smiling more and more. »I brought it to thee a bit hungry because I knowed tha'd like to see it feed.«

He knelt down by the sofa and took a feeding-bottle from his pocket.

»Come on, little 'un,« he said, turning the small woolly white head with a gentle brown hand. »This is what tha's after. Tha'll get more out o' this than tha' will out o' silk velvet coats. There now,« and he pushed the rubber tip of the bottle into the nuzzling mouth and the lamb began to suck it with ravenous ecstasy.

After that there was no wondering what to say. By the time the lamb fell asleep questions poured forth and Dickon answered them all. He told them how he had found the lamb just as the sun was rising three mornings ago. He had been standing on the moor listening to a skylark and watching him swing higher and higher into the sky until he was only a speck in the heights of blue.

»I'd almost lost him but for his song an' I was wonderin' how a chap could hear it when it seemed as if he'd get out o' th' world in a minute—an' just then I heard somethin' else far off among th' gorse bushes. It was a weak bleatin' an' I knowed it was a new lamb as was hungry an' I knowed it wouldn't be hungry if it hadn't lost its mother somehow, so I set off searchin'. Eh! I did have a look for it. I went in an' out among th' gorse bushes an' round an' round an' I always seemed to take th' wrong turnin'. But at last I seed a bit o' white by a rock on top o' th' moor an' I climbed up an' found th' little 'un half dead wi' cold an' clemmin'.«

While he talked, Soot flew solemnly in and out of the open window and cawed remarks about the scenery while Nut and Shell made excursions into the big trees outside and ran up and down trunks and explored branches. Captain curled up near Dickon, who sat on the hearth-rug from preference.

They looked at the pictures in the gardening books and Dickon knew all the flowers by their country names and knew exactly which ones were already growing in the secret garden.

»I couldna' say that there name,« he said, pointing to one under which was written »Aquilegia,« »but us calls that a columbine, an' that there one it's a snapdragon and they both grow wild in hedges, but these is garden ones an' they're bigger an' grander. There's some big clumps o' columbine in th' garden. They'll look like a bed o' blue an' white butterflies flutterin' when they're out.«

»I'm going to see them,« cried Colin. »I am going to see them!«

»Aye, that tha' mun,« said Mary quite seriously. »An tha' munnot lose no time about it.«